\documentclass[conference]{IEEEtran}
\IEEEoverridecommandlockouts

% -------------------------------------------------------
% Packages
% -------------------------------------------------------
\usepackage{cite}
\usepackage{amsmath,amssymb,amsfonts}
\usepackage{algorithmic}
\usepackage{graphicx}
\usepackage{textcomp}
\usepackage{xcolor}
\usepackage{booktabs}
\usepackage{array}
\usepackage{url}
\usepackage{hyperref}
\usepackage{multirow}
\usepackage{tabularx}
\usepackage[T1]{fontenc}

\hypersetup{
    colorlinks=true,
    linkcolor=blue,
    citecolor=blue,
    urlcolor=blue
}

% -------------------------------------------------------
% Title & Authors
% -------------------------------------------------------
\begin{document}


\title{SynthInspect: A Conditional Generative Augmentation Pipeline\\
for Rare Steel Surface Defect Detection\\}


\author{
\IEEEauthorblockN{Mohit Goyal}
\IEEEauthorblockA{
\textit{Department of Computer Science and Engineering} \\
\textit{Bharati Vidyapeeth's College of Engineering} \\
New Delhi, India \\
mohitgoyal09042006@gmail.com}
\and
\IEEEauthorblockN{Angel Gupta}
\IEEEauthorblockA{
\textit{Department of Computer Science and Engineering} \\
\textit{Bharati Vidyapeeth's College of Engineering} \\
New Delhi, India \\
angelgupt16@gmail.com}
}



\maketitle

% -------------------------------------------------------
% ABSTRACT
% -------------------------------------------------------
\begin{abstract}
Industrial steel surface inspection is critically constrained by scarce
and imbalanced annotated defect samples. In the Severstal Steel Defect
Detection benchmark, 63\% of images are defect-free while the rarest
defect class contains only approximately 320 annotated instances,
making robust supervised training difficult and annotation expansion
prohibitively expensive. This project proposes \textit{SynthInspect},
a conditional generative augmentation pipeline that learns defect
distributions directly from Severstal line-scan images to
systematically improve rare-defect detection.

We design a DG-GAN architecture with a U-Net-inspired generator
incorporating Atrous Spatial Pyramid Pooling (ASPP) and defect-aware
attention mechanisms for 1600$\times$256 wide-aspect images, coupled
with a multi-scale PatchGAN discriminator trained via WGAN-GP. The
generator is conditioned on defect class, spatial location, size, and
severity, enabling controllable synthesis of diverse defects on
realistic steel backgrounds. Generated samples undergo
discriminator-score quality filtering and are integrated into a
curriculum-based training loop with dynamically scheduled real:synthetic
mixing ratios. We hypothesize a 15--25\% relative improvement in
rare-class average precision versus traditional augmentation alone,
alongside a 50--90\% reduction in real annotation requirements.
Validation comprises ablation studies on Severstal and cross-dataset
transfer to NEU Surface Defect, with a fully open-source PyTorch
implementation targeting IECON 2026 submission.
\end{abstract}

\begin{IEEEkeywords}
Generative Adversarial Networks, Steel Defect Detection, Data Augmentation,
Curriculum Learning, WGAN-GP, Severstal, Image Synthesis
\end{IEEEkeywords}

% -------------------------------------------------------
% SECTION I: INTRODUCTION
% -------------------------------------------------------
\section{Introduction}

Automated visual inspection of steel surfaces is a critical quality
control step in modern manufacturing. However, supervised deep learning
approaches for defect detection are limited by severe data imbalance:
the majority of production images show defect-free surfaces, while
rare but high-impact defect types have only a handful of annotated
examples. In the widely-used Severstal Steel Defect Detection
benchmark \cite{severstal_kaggle}, Class~2 (pitted surface) contains
approximately 320 instances against over 7,000 defect-free images,
a class imbalance ratio of roughly 45:1.

Traditional countermeasures---geometric and photometric
augmentations---operate on a fixed set of defect templates and cannot
synthesize genuinely new defect morphologies or placements
\cite{severstal_aug_study}. This motivates our use of conditional
Generative Adversarial Networks (GANs) to \emph{learn} the defect
distribution from real Severstal imagery and generate novel,
mask-paired defect examples on demand.

% -------------------------------------------------------
% SECTION II: LITERATURE REVIEW
% -------------------------------------------------------
\section{Literature Review}

\subsection{Traditional Augmentation and Procedural Synthesis}

Early work on steel defect recognition relied on geometric and
photometric augmentations (rotation, flipping, brightness jittering)
and procedural simulation. Boikov et al.\ \cite{boikov2021synthetic}
generated fully synthetic steel surface images via procedural rendering
and trained CNNs for detection and classification on Severstal-like
conditions. Despite strong reported performance, these synthetic textures
are not learned from actual Severstal imagery and lack fine-grained mask
conditioning over defect location and severity. Consequently, such methods
cannot address the long-tail imbalance problem for specific rare-class
instance generation. A systematic study of augmentation strategies on
Severstal U-Net baselines \cite{severstal_aug_study} confirms that
naive augmentation alone yields limited gains on this dataset.

\subsection{GAN-Based Defect Synthesis}

GAN-based approaches have demonstrated significant potential for
industrial defect augmentation. Fusible-feature GANs for steel strip
defects under few-shot constraints \cite{fusible_gan} fuse background
and defect features to expand limited datasets and improve recognition
in highly imbalanced settings. DG-GAN variants \cite{dg_gan_surface}
employ U-shaped generators with spatially adaptive normalization and
custom perceptual losses for high-quality synthesis, reporting improved
FID and downstream classification improvements. A sketch-guided,
SPADE-like architecture \cite{sketch_spade_steel} enables controllable
generation via edge and texture conditioning for steel strip surfaces.
However, these methods typically: (a) evaluate generation quality in
isolation; (b) use fixed synthetic-to-real ratios without adaptive
scheduling; and (c) omit systematic per-class AP evaluation on canonical
benchmarks such as Severstal.

\subsection{Diffusion Models for Defect Data}

Recent work fine-tunes latent diffusion models on steel datasets for
segmentation enhancement \cite{diffusion_steel_seg}. While achieving
high fidelity, diffusion-based approaches require substantially more
compute and more complex conditional control mechanisms, making them
less practical for rapid deployment in resource-constrained industrial
settings within short research timelines.

\subsection{Research Gap}

Existing approaches exhibit three principal limitations:
\begin{enumerate}
    \item Absence of fine-grained class, location, size, and severity
          conditioning directly trained on Severstal line-scan images.
    \item No closed-loop curriculum that dynamically adjusts the
          real:synthetic ratio based on downstream validation
          performance.
    \item Insufficient systematic evaluation of rare-class detection
          gains and annotation reduction on a widely-benchmarked
          industrial dataset.
\end{enumerate}
Our work addresses all three gaps through a Severstal-specific,
mask-conditioned DG-GAN with discriminator-score quality filtering and
a validation-driven curriculum scheduler---detailed in
Section~\ref{sec:methodology}.

% -------------------------------------------------------
% SECTION III: PROPOSED METHODOLOGY
% -------------------------------------------------------
\section{Proposed Methodology}
\label{sec:methodology}

Fig.~\ref{fig:pipeline} provides a full overview of the
\textit{SynthInspect} pipeline. The system has three integrated
stages: (1) conditional defect generation, (2) quality-filtered
curriculum integration, and (3) downstream detection with systematic
evaluation.

\begin{figure*}[t]
    \centering
    \includegraphics[width=\textwidth]{pipeline_image.png}
    \caption{Overview of the SynthInspect closed-loop generative
    augmentation pipeline. Severstal raw images and RLE masks are
    preprocessed and split into two branches: DG-GAN training and
    baseline detector training. The conditional generator and
    multi-scale PatchGAN discriminator learn to synthesize realistic,
    mask-conditioned defect patches. Synthetic samples are filtered by
    discriminator score, combined with real data, and fed to a
    curriculum scheduler that dynamically adjusts the real:synthetic
    ratio using validation mAP. The final detector is evaluated on a
    held-out real test set.}
    \label{fig:pipeline}
\end{figure*}

\subsection{Stage I: Preprocessing and Data Preparation}

Raw Severstal images (1600$\times$256 px) and run-length encoded (RLE)
masks from \texttt{train.csv} are decoded into binary segmentation
masks. The dataset is split in a stratified manner (80\% train,
10\% validation, 10\% test) to preserve rare-class representation.
Images are resized or cropped to 512$\times$512 inputs with
aspect-preserving padding, followed by ImageNet-mean normalization.
Standard geometric augmentations (horizontal/vertical flip,
random rotation $\pm$15$^\circ$, brightness/contrast jitter) are
applied during baseline training only.

\subsection{Stage II: Conditional DG-GAN}

\subsubsection{Generator Architecture}
The generator $G$ is a U-Net-style encoder-decoder with five
downsampling blocks using residual connections, producing feature maps
at multiple spatial resolutions. Atrous Spatial Pyramid Pooling (ASPP)
with dilation rates $\{6, 12, 18\}$ is inserted at the bottleneck
to capture both fine-grained textures and large-scale defect structures.
Defect-aware attention (spatial + channel + class-conditional) is
applied in decoder layers to localize synthesis to the specified region.

The conditioning vector $\mathbf{c} = [c, \mathbf{b}, s, v]$ encodes:
\begin{itemize}
    \item $c \in \{1,2,3,4\}$: defect class (one-hot),
    \item $\mathbf{b} = (x,y,w,h)$: bounding box location,
    \item $s \in [0,1]$: normalized defect area,
    \item $v \in [0,1]$: severity estimate.
\end{itemize}
$\mathbf{c}$ is injected via class-conditional batch normalization and
concatenated with spatial feature maps after projection.

\subsubsection{Discriminator Architecture}
The discriminator $D$ is a multi-scale PatchGAN operating at three
resolutions (full, half, quarter) using shared parameters via strided
average pooling. Each scale outputs a patch-level realism map, and all
maps are averaged for the final score.

\subsubsection{Training Objective}
The training loss combines:
\begin{align}
    \mathcal{L}_{\text{WGAN-GP}} &= \mathbb{E}[D(\tilde{x})] -
    \mathbb{E}[D(x)] + \lambda_{\text{GP}}\,\hat{\mathcal{L}}_{\text{GP}},
    \label{eq:wgan} \\
    \mathcal{L}_G &= -\mathbb{E}[D(\tilde{x})] + \lambda_1\mathcal{L}_1
    + \lambda_p\mathcal{L}_{\text{perc}},
    \label{eq:gen_loss}
\end{align}
where $\lambda_{\text{GP}}=10$, $\lambda_1=10$, and
$\lambda_p=0.1$. $\mathcal{L}_1$ is an L1 reconstruction loss against
the real defect-region content, and $\mathcal{L}_{\text{perc}}$ is a
VGG-19 perceptual loss on feature activations from layers
\texttt{relu2\_2} and \texttt{relu4\_4}.

\subsection{Stage III: Quality Filtering and Curriculum Integration}

\subsubsection{Discriminator-Score Quality Filter}
After training, $G$ is used to generate a pool of $N_{\text{pool}}$
synthetic images. The discriminator assigns each generated sample a
realism score $\sigma_i = D(G(z_i, \mathbf{c}_i))$. Samples are
retained if:
\begin{equation}
    \sigma_i \in [q_{30}, q_{70}],
    \label{eq:filter}
\end{equation}
where $q_{30}$ and $q_{70}$ are the 30th and 70th percentiles of
$\{\sigma_i\}$. This retains samples that are realistic yet
challenging, discarding obvious failures (low $\sigma$) and trivially
easy samples (very high $\sigma$, near-memorization).

\subsubsection{Curriculum Scheduler}
The training dataset at epoch $t$ is constructed as:
\begin{equation}
    \mathcal{D}_t = (1-\rho_t)\,\mathcal{D}_{\text{real}} +
    \rho_t\,\mathcal{D}_{\text{syn}},
    \label{eq:mix}
\end{equation}
with the mixing ratio $\rho_t$ initialized to 0.10 and updated every
$\Delta$ epochs according to:
\begin{equation}
    \rho_{t+\Delta} =
    \begin{cases}
        \min(\rho_t + 0.05,\ 0.50) & \text{if } \Delta\text{mAP}_{\text{val}} > 0 \\
        \max(\rho_t - 0.05,\ 0.10) & \text{if } \Delta\text{mAP}_{\text{val}} < -\epsilon \\
        \rho_t & \text{otherwise,}
    \end{cases}
    \label{eq:curriculum}
\end{equation}
where $\epsilon = 0.005$ is a tolerance threshold for mAP degradation.

\subsection{Stage IV: Downstream Detection}

Two detector architectures are evaluated:
\begin{itemize}
    \item \textbf{U-Net with ResNet-34 encoder}: segmentation model
          from \texttt{segmentation\_models.pytorch} \cite{khornlund_repo},
          trained with BCE + Dice loss.
    \item \textbf{YOLOv8-m}: bounding-box detection for architecture
          agnosticism evaluation.
\end{itemize}
Both receive the curriculum-mixed dataset $\mathcal{D}_t$ and are
evaluated on a held-out real test set (no synthetic images) to avoid
distribution leakage.

% -------------------------------------------------------
% SECTION IV: DATASET AND EXPERIMENT PLAN
% -------------------------------------------------------
\section{Dataset and Experiment Plan}

\subsection{Datasets}

\subsubsection{Primary: Severstal Steel Defect Detection}
Table~\ref{tab:dataset} summarizes the dataset characteristics.

\begin{table}[h]
\centering
\caption{Severstal Steel Defect Detection Dataset Summary}
\label{tab:dataset}
\begin{tabular}{ll}
\toprule
\textbf{Attribute} & \textbf{Specification} \\
\midrule
Source          & Kaggle Competition, 2019 \cite{severstal_kaggle} \\
Total images    & 18,074 (12,568 train / 5,506 test) \\
Resolution      & 1600$\times$256 px (line-scan) \\
Annotation      & Pixel-level RLE masks \\
Classes         & 4 defect types + background \\
Defect-free     & $\approx$63\% of all images \\
Rarest class    & Class 2: $\approx$320 instances (45:1 ratio) \\
Split           & 80 / 10 / 10 (stratified by class) \\
\bottomrule
\end{tabular}
\end{table}

\textbf{Implementation references:}
We build upon established PyTorch repositories: khornlund
\cite{khornlund_repo} for the SMP-based U-Net baseline, and Diyago
\cite{diyago_repo} for Albumentations-based augmentation pipelines
and attention blocks (Top-2\% competition solution). EDA and RLE
decoding follow Aleksandra's Kaggle notebook \cite{aleksandra_eda}.

\subsubsection{Secondary: NEU Surface Defect}
1,800 grayscale 200$\times$200 images across six defect categories
(Northeastern University, China). Used exclusively for cross-dataset
transfer evaluation to assess generalization.

\subsection{Planned Experiments}

Table~\ref{tab:experiments} lists the five core ablation experiments.

\begin{table}[h]
\centering
\caption{Planned Experiments and Hypotheses}
\label{tab:experiments}
\renewcommand{\arraystretch}{1.25}
\begin{tabularx}{\columnwidth}{>{\bfseries}l X X}
\toprule
\textbf{Experiment} & \textbf{Variants} & \textbf{Hypothesis} \\
\midrule
E1: Baseline
& Real-only; Real + traditional aug
& Traditional aug $<$5\% mAP gain \\
\addlinespace
E2: Ratio ablation
& $\rho \in \{0, 0.10, 0.25, 0.50, 0.75\}$
& Optimal $\rho \approx 0.30$--$0.50$ \\
\addlinespace
E3: Curriculum vs.\ fixed
& Scheduled $\rho$ (10\%$\to$50\%) vs.\ best fixed $\rho$
& Curriculum outperforms by $>$5\% mAP \\
\addlinespace
E4: Filtering ablation
& With vs.\ without D-score filter
& Filtering improves sample efficiency \\
\addlinespace
E5: Annotation reduction
& Remove 50\%, 75\%, 90\% rare-class labels
& $<$3\% mAP drop with 50\% removal \\
\bottomrule
\end{tabularx}
\end{table}

\subsection{Evaluation Metrics}

Table~\ref{tab:metrics} defines all metrics used for evaluation.

\begin{table}[h]
\centering
\caption{Evaluation Metrics and Targets}
\label{tab:metrics}
\renewcommand{\arraystretch}{1.2}
\begin{tabularx}{\columnwidth}{l X l}
\toprule
\textbf{Category} & \textbf{Metric} & \textbf{Target} \\
\midrule
Generation quality  & FID (lower is better)         & $<50$ \\
Perceptual quality  & LPIPS (lower is better)       & $<0.30$ \\
Detection           & mAP@0.50                      & $>0.75$ \\
Detection           & mAP@0.50:0.95                 & $>0.50$ \\
Rare-class focus    & Per-class AP (Class~2)        & $>0.60$ \\
Segmentation        & Mean IoU / Dice               & $>0.70$ / $>0.80$ \\
Annotation saving   & Effective label reduction     & 50--90\% \\
\bottomrule
\end{tabularx}
\end{table}

\subsection{Implementation Timeline}

Table~\ref{tab:timeline} shows the week-by-week implementation plan
aligned with the 10-week project schedule.

\begin{table}[h]
\centering
\caption{Implementation Roadmap (10-Week Schedule)}
\label{tab:timeline}
\renewcommand{\arraystretch}{1.2}
\begin{tabularx}{\columnwidth}{c X l}
\toprule
\textbf{Week} & \textbf{Task} & \textbf{Deliverable} \\
\midrule
1   & EDA, RLE decoding, data loading,
      clone khornlund baseline \cite{khornlund_repo}
    & Working data loader; class distribution plots \\
2   & Albumentations pipeline \cite{diyago_repo};
      attention blocks; strong baseline
    & Baseline mAP on real data only \\
3   & DG-GAN generator (U-Net + ASPP)
    & 512$\times$512 defect generation \\
4   & Multi-scale PatchGAN + WGAN-GP;
      stable GAN checkpoint
    & FID $<$ 150 on defect patches \\
5   & Quality filtering; 10K synthetic dataset
    & FID $<$ 100; RLE masks saved \\
6   & Curriculum integration;
      mixed training experiments
    & mAP vs.\ baseline; E1, E2 results \\
7   & Curriculum vs.\ fixed (E3);
      filtering ablation (E4)
    & Full ablation tables \\
8   & Annotation reduction (E5);
      cross-dataset on NEU
    & Transfer and annotation results \\
9   & Paper drafting; figures; code cleanup
    & Draft IEEE paper \\
10  & Revision; open-source release; submission
    & IECON 2026 submission \\
\bottomrule
\end{tabularx}
\end{table}

% -------------------------------------------------------
% Risk Mitigation
% -------------------------------------------------------
\section{Risk Mitigation}

To address potential technical risks, we implement the following
contingency strategies:

\begin{itemize}
    \item \textbf{GAN training instability:} Checkpoint every epoch
          with FID evaluation. If conditional DG-GAN fails to converge
          by Week~4, fallback to (a) unconditional GAN on defect patches
          with post-hoc classification, or (b) fine-tuned Stable
          Diffusion with LoRA.
    \item \textbf{Compute constraints:} Pre-allocated cloud GPU
          instances (AWS/GCP) with mixed-precision training and
          gradient checkpointing.
    \item \textbf{Dataset access:} Confirmed Kaggle availability;
          NEU dataset as verified backup.
\end{itemize}

% -------------------------------------------------------
% SECTION V: CONCLUSION
% -------------------------------------------------------
\section{Conclusion}

This milestone report presents \textit{SynthInspect}, a conditional
generative augmentation framework for steel surface defect detection.
The proposed system addresses critical limitations of prior work by
integrating a Severstal-specific mask-conditioned DG-GAN with
discriminator-score quality filtering and a validation-driven
curriculum scheduler. The methodology diagram (Fig.~\ref{fig:pipeline})
and experiment plan establish a rigorous evaluation framework targeting
15--25\% rare-class AP improvement with 50--90\% annotation savings.
Upcoming work will focus on delivering a working PyTorch baseline
(Week~1--2) and a stable conditional generator (Week~3--4) as the
foundation for all downstream experiments.

% -------------------------------------------------------
% REFERENCES
% -------------------------------------------------------
\begin{thebibliography}{00}

\bibitem{severstal_kaggle}
Severstal and Kaggle, ``Severstal: Steel Defect Detection,''
Kaggle Competition Dataset, 2019. [Online]. Available:
\url{https://www.kaggle.com/c/severstal-steel-defect-detection}

\bibitem{severstal_aug_study}
Anonymous, ``A Systematic Augmentation Study for Steel Defect
Segmentation on Severstal,'' in \textit{Proc. ACE}, 2023.
[Online]. Available:
\url{https://www.ewadirect.com/proceedings/ace/article/view/23131}

\bibitem{boikov2021synthetic}
V.~Boikov, V.~Payor, R.~Savelev, and A.~Kolesnikov,
``Synthetic Data Generation for Steel Defect Detection and
Classification Using Deep Learning,'' \textit{Symmetry}, vol.~13,
no.~7, p.~1176, 2021.

\bibitem{fusible_gan}
Y.~Liu et al., ``Few-Shot Steel Surface Defect Detection via
Fusible-Feature GAN for Image Synthesis,'' \textit{IEEE Access},
vol.~11, pp.~30,224--30,237, 2023. [Online]. Available:
\url{https://pmc.ncbi.nlm.nih.gov/articles/PMC10054326/}

\bibitem{dg_gan_surface}
X.~Chen et al., ``High-Quality Defect Generation for Industrial
Surface Inspection Using Generative Adversarial Networks,''
\textit{Sensors}, vol.~23, no.~6, p.~3216, 2023. [Online]. Available:
\url{https://www.mdpi.com/1424-8220/23/6/3216}

\bibitem{sketch_spade_steel}
Z.~Li et al., ``Sketch-Guided Spatially Adaptive Defect Synthesis
for Steel Strip Surfaces,'' \textit{Measurement Science and
Technology}, vol.~35, no.~5, 2024. DOI:
\url{https://iopscience.iop.org/article/10.1088/1361-6501/ad1eb6}

\bibitem{diffusion_steel_seg}
H.~Zhang et al., ``Enhancing Steel Surface Defect Segmentation via
Latent Diffusion Model Augmentation,'' \textit{Sensors}, vol.~23,
no.~19, 2023. [Online]. Available:
\url{https://pmc.ncbi.nlm.nih.gov/articles/PMC11436218/}

\bibitem{khornlund_repo}
K.~Hornlund, ``Severstal Steel Defect Detection,'' GitHub, 2019.
[Online]. Available:
\url{https://github.com/khornlund/severstal-steel-defect-detection}

\bibitem{diyago_repo}
D.~Ignatov, ``Severstal-Steel-Defect-Detection (Top-2\%
Solution),'' GitHub, 2019. [Online]. Available:
\url{https://github.com/Diyago/Severstal-Steel-Defect-Detection}

\bibitem{aleksandra_eda}
A.~Deis, ``Steel Defect Detection EDA,'' Kaggle Notebook,
2019. [Online]. Available:
\url{https://www.kaggle.com/code/aleksandradeis/steel-defect-detection-eda}

\end{thebibliography}

\end{document}